% GENERATED FILE. Edit canonical lessons, then run npm run generate:books.
% canonical-source-hash: fnv1a64:3e47ee75ad9350c2
% canonical-lessons: ML-C550-jolisthalam, ML-C550-melnottam, ML-C550-ekopippikkuka, ML-C550-kutiyalocikkuka, ML-C550-itavela

\chapter{Managing and Consulting}
\label{ch:ml-managing-and-consulting}
\hlchaptermodality{550}

\begin{chapteropening}
\textbf{By the end of this chapter:} \emph{I can say a workplace, supervision, to coordinate, to consult and an interval.}

Complete the last lesson of chapter 550: I can say a workplace, supervision, to coordinate, to consult and an interval.
\end{chapteropening}

\section[\texorpdfstring{\ml{ജോലിസ്ഥലം} (jōlisthalaṁ)}{ജോലിസ്ഥലം (jōlisthalaṁ)}]{\ml{ജോലിസ്ഥലം} (jōlisthalaṁ) — a workplace}

\label{lesson:ML-C550-jolisthalam}



\begin{quote}
Before the new one: say the Malayalam for a pay rise, then the Malayalam for to bear a load.
\end{quote}

\subsection*{You'll want to know: \ml{ജോലിസ്ഥലം}}
\textbf{\ml{ജോലിസ്ഥലം}} — \emph{jōlisthalaṁ} — \textquotedblleft{}a workplace\textquotedblright{}.

\begin{quote}
\ml{അവളുടെ} \ml{ജോലിസ്ഥലം} \ml{വീട്ടിൽ} \ml{നിന്ന്} \ml{ദൂരെയാണ്}.
\end{quote}

\begin{quote}
\emph{avaḷuṭe jōlisthalaṁ vīṭṭil ninnŭ dūreyāṇŭ.}
\end{quote}

\begin{quote}
Her workplace is far from home.
\end{quote}

\ml{ജോലി} work + \ml{സ്ഥലം} a place.

\begin{cousinweb}
\textbf{\ml{ജോലിസ്ഥലം}} (\emph{jōlisthalaṁ}) is a Malayalam compound built from \ml{ജോലി}, \textquotedblleft{}work\textquotedblright{}, and \ml{സ്ഥലം}, from Sanskrit \emph{sthala}, \textquotedblleft{}ground, a place\textquotedblright{} (CDIAL 13744).
\end{cousinweb}

\begin{grammarlens}[title={What you've built}]
One more word for describing what you lived and felt.
\end{grammarlens}

\subsection*{Your turn}
\begin{itemize}
\raggedright
  \item \emph{Say it:} \emph{jōlisthalaṁ}
  \item \emph{Say it:} \emph{jōlisthalaṁ}, once more
  \item \emph{Say it:} \emph{avaḷuṭe jōlisthalaṁ vīṭṭil ninnŭ dūreyāṇŭ.} — Her workplace is far from home.
  \item \emph{Recall:} say the Malayalam for a pay rise, then the Malayalam for to bear a load, then say \emph{jōlisthalaṁ} again
\end{itemize}

\begin{tcolorbox}[breakable,colback=teal!4,colframe=teal!35!black,title={Before you move on}]
What does \ml{ജോലിസ്ഥലം} mean? (\textquotedblleft{}A workplace\textquotedblright{}.) Say it once more.
\end{tcolorbox}
\section[\texorpdfstring{\ml{മേൽനോട്ടം} (mēlnōṭṭaṁ)}{മേൽനോട്ടം (mēlnōṭṭaṁ)}]{\ml{മേൽനോട്ടം} (mēlnōṭṭaṁ) — supervision}

\label{lesson:ML-C550-melnottam}



\begin{quote}
Before the new one: say the Malayalam for to bear a load, then the Malayalam for a workplace.
\end{quote}

\subsection*{You'll want to know: \ml{മേൽനോട്ടം}}
\textbf{\ml{മേൽനോട്ടം}} — \emph{mēlnōṭṭaṁ} — \textquotedblleft{}supervision\textquotedblright{}.

\begin{quote}
\ml{പണിയുടെ} \ml{മേൽനോട്ടം} \ml{മേസ്തിരിക്കായിരുന്നു}.
\end{quote}

\begin{quote}
\emph{paṇiyuṭe mēlnōṭṭaṁ mēstirikkāyirunnu.}
\end{quote}

\begin{quote}
The mason was in charge of supervising the work.
\end{quote}

\ml{മേൽ} above + \ml{നോട്ടം} a look: watching from above.

\begin{cousinweb}
\textbf{\ml{മേൽനോട്ടം}} (\emph{mēlnōṭṭaṁ}) is native Dravidian, a compound of \ml{മേൽ}, \textquotedblleft{}above\textquotedblright{} (DEDR 5086), and \ml{നോട്ടം}, \textquotedblleft{}a look, an examination\textquotedblright{} (DEDR 3794).
\end{cousinweb}

\begin{grammarlens}[title={What you've built}]
One more word for describing what you lived and felt.
\end{grammarlens}

\subsection*{Your turn}
\begin{itemize}
\raggedright
  \item \emph{Say it:} \emph{mēlnōṭṭaṁ}
  \item \emph{Say it:} \emph{mēlnōṭṭaṁ}, once more
  \item \emph{Say it:} \emph{paṇiyuṭe mēlnōṭṭaṁ mēstirikkāyirunnu.} — The mason was in charge of supervising the work.
  \item \emph{Recall:} say the Malayalam for to bear a load, then the Malayalam for a workplace, then say \emph{mēlnōṭṭaṁ} again
\end{itemize}

\begin{tcolorbox}[breakable,colback=teal!4,colframe=teal!35!black,title={Before you move on}]
What does \ml{മേൽനോട്ടം} mean? (\textquotedblleft{}Supervision\textquotedblright{}.) Say it once more.
\end{tcolorbox}
\section[\texorpdfstring{\ml{ഏകോപിപ്പിക്കുക} (ēkōpippikkuka)}{ഏകോപിപ്പിക്കുക (ēkōpippikkuka)}]{\ml{ഏകോപിപ്പിക്കുക} (ēkōpippikkuka) — to coordinate}

\label{lesson:ML-C550-ekopippikkuka}



\begin{quote}
Before the new one: say the Malayalam for a workplace, then the Malayalam for supervision.
\end{quote}

\subsection*{You'll want to know: \ml{ഏകോപിപ്പിക്കുക}}
\textbf{\ml{ഏകോപിപ്പിക്കുക}} — \emph{ēkōpippikkuka} — \textquotedblleft{}to coordinate\textquotedblright{}.

\begin{quote}
\ml{പരിപാടിയുടെ} \ml{എല്ലാ} \ml{ജോലികളും} \ml{അവൾ} \ml{ഏകോപിപ്പിച്ചു}.
\end{quote}

\begin{quote}
\emph{paripāṭiyuṭe ellā jōlikaḷuṁ avaḷ ēkōpippiccu.}
\end{quote}

\begin{quote}
She coordinated all the work for the programme.
\end{quote}

Past: \ml{ഏകോപിപ്പിച്ചു}, \emph{ēkōpippiccu}.

The noun is \ml{ഏകോപനം}, coordination: making many into one.

\begin{cousinweb}
\textbf{\ml{ഏകോപിപ്പിക്കുക}} (\emph{ēkōpippikkuka}) is a learned verb built on Sanskrit \emph{eka}, \textquotedblleft{}one\textquotedblright{} (Monier-Williams), with the causative ending -\ml{ിപ്പിക്കുക}; the exact model of the formation is not settled.
\end{cousinweb}

\begin{grammarlens}[title={What you've built}]
One more word for describing what you lived and felt.
\end{grammarlens}

\subsection*{Your turn}
\begin{itemize}
\raggedright
  \item \emph{Say it:} \emph{ēkōpippikkuka}
  \item \emph{Say it:} \emph{ēkōpippikkuka}, once more
  \item \emph{Say it:} \emph{paripāṭiyuṭe ellā jōlikaḷuṁ avaḷ ēkōpippiccu.} — She coordinated all the work for the programme.
  \item \emph{Recall:} say the Malayalam for a workplace, then the Malayalam for supervision, then say \emph{ēkōpippikkuka} again
\end{itemize}

\begin{tcolorbox}[breakable,colback=teal!4,colframe=teal!35!black,title={Before you move on}]
What does \ml{ഏകോപിപ്പിക്കുക} mean? (\textquotedblleft{}To coordinate\textquotedblright{}.) Say it once more.
\end{tcolorbox}
\section[\texorpdfstring{\ml{കൂടിയാലോചിക്കുക} (kūṭiyālōcikkuka)}{കൂടിയാലോചിക്കുക (kūṭiyālōcikkuka)}]{\ml{കൂടിയാലോചിക്കുക} (kūṭiyālōcikkuka) — to consult; to confer}

\label{lesson:ML-C550-kutiyalocikkuka}



\begin{quote}
Before the new one: say the Malayalam for supervision, then the Malayalam for to coordinate.
\end{quote}

\subsection*{You'll want to know: \ml{കൂടിയാലോചിക്കുക}}
\textbf{\ml{കൂടിയാലോചിക്കുക}} — \emph{kūṭiyālōcikkuka} — \textquotedblleft{}to consult; to confer\textquotedblright{}.

\begin{quote}
\ml{തീരുമാനം} \ml{എടുക്കും} \ml{മുൻപ്} \ml{മാനേജർ} \ml{ടീമുമായി} \ml{കൂടിയാലോചിച്ചു}.
\end{quote}

\begin{quote}
\emph{tīrumānaṁ eṭukkuṁ munpŭ mānējar ṭīmumāyi kūṭiyālōciccu.}
\end{quote}

\begin{quote}
Before deciding, the manager consulted the team.
\end{quote}

Past: \ml{കൂടിയാലോചിച്ചു}, \emph{kūṭiyālōciccu}.

\ml{കൂടി} together + \ml{ആലോചിക്കുക} to think over. \ml{കൂടിയാലോചന} is the noun.

\begin{cousinweb}
\textbf{\ml{കൂടിയാലോചിക്കുക}} (\emph{kūṭiyālōcikkuka}) is a Malayalam compound built from \ml{കൂടി}, \textquotedblleft{}together\textquotedblright{}, which is native (DEDR 1882), and \ml{ആലോചിക്കുക}, from Sanskrit \emph{ālocana}, \textquotedblleft{}considering, reflection\textquotedblright{} (Monier-Williams).
\end{cousinweb}

\begin{grammarlens}[title={What you've built}]
One more word for describing what you lived and felt.
\end{grammarlens}

\subsection*{Your turn}
\begin{itemize}
\raggedright
  \item \emph{Say it:} \emph{kūṭiyālōcikkuka}
  \item \emph{Say it:} \emph{kūṭiyālōcikkuka}, once more
  \item \emph{Say it:} \emph{tīrumānaṁ eṭukkuṁ munpŭ mānējar ṭīmumāyi kūṭiyālōciccu.} — Before deciding, the manager consulted the team.
  \item \emph{Recall:} say the Malayalam for supervision, then the Malayalam for to coordinate, then say \emph{kūṭiyālōcikkuka} again
\end{itemize}

\begin{tcolorbox}[breakable,colback=teal!4,colframe=teal!35!black,title={Before you move on}]
What does \ml{കൂടിയാലോചിക്കുക} mean? (\textquotedblleft{}To consult; to confer\textquotedblright{}.) Say it once more.
\end{tcolorbox}
\section[\texorpdfstring{\ml{ഇടവേള} (iṭavēḷa)}{ഇടവേള (iṭavēḷa)}]{\ml{ഇടവേള} (iṭavēḷa) — an interval; a break}

\label{lesson:ML-C550-itavela}



\begin{quote}
Before the new one: say the Malayalam for to coordinate, then the Malayalam for to consult.
\end{quote}

\subsection*{You'll want to know: \ml{ഇടവേള}}
\textbf{\ml{ഇടവേള}} — \emph{iṭavēḷa} — \textquotedblleft{}an interval; a break\textquotedblright{}.

\begin{quote}
\ml{രണ്ട്} \ml{മണിക്കൂർ} \ml{ജോലിക്ക്} \ml{ശേഷം} \ml{പത്ത്} \ml{മിനിറ്റ്} \ml{ഇടവേള} \ml{ഉണ്ട്}.
\end{quote}

\begin{quote}
\emph{raṇṭŭ maṇikkūr jōlikkŭ śēṣaṁ pattŭ miniṟṟŭ iṭavēḷa uṇṭŭ.}
\end{quote}

\begin{quote}
After two hours of work there is a ten-minute break.
\end{quote}

\ml{ഇട} a gap + \ml{വേള} a time; at the cinema it is the interval.

\begin{cousinweb}
\textbf{\ml{ഇടവേള}} (\emph{iṭavēḷa}) is a Malayalam compound built from \ml{ഇട}, \textquotedblleft{}an interval\textquotedblright{}, which is native (DEDR 448), and \ml{വേള}, from Sanskrit \emph{velā}, \textquotedblleft{}a limit, a boundary\textquotedblright{} and so a point of time (Monier-Williams).
\end{cousinweb}

\begin{grammarlens}[title={What you've built}]
One more word for describing what you lived and felt.
\end{grammarlens}

\subsection*{Your turn}
\begin{itemize}
\raggedright
  \item \emph{Say it:} \emph{iṭavēḷa}
  \item \emph{Say it:} \emph{iṭavēḷa}, once more
  \item \emph{Say it:} \emph{raṇṭŭ maṇikkūr jōlikkŭ śēṣaṁ pattŭ miniṟṟŭ iṭavēḷa uṇṭŭ.} — After two hours of work there is a ten-minute break.
  \item \emph{Recall:} say the Malayalam for to coordinate, then the Malayalam for to consult, then say \emph{iṭavēḷa} again
\end{itemize}

\begin{tcolorbox}[breakable,colback=teal!4,colframe=teal!35!black,title={Before you move on}]
What does \ml{ഇടവേള} mean? (\textquotedblleft{}An interval; a break\textquotedblright{}.) Say it once more.
\end{tcolorbox}
